\documentclass[10pt,a4paper]{amsart}
\usepackage[margin=19mm]{geometry}
\usepackage{amsmath,amssymb,mathtools}
\usepackage[scr=rsfs,cal=euler]{mathalfa}
\usepackage{xfrac}
\usepackage[unicode,hidelinks,bookmarks=false]{hyperref}
\newcommand{\ie}{i.e.\ }
\newcommand{\cf}{cf.\ }
\newcommand{\resp}{respectively\ }
\newcommand{\Subsection}{\subsection}
\providecommand{\nobreakdash}{\nobreak\hbox{-}}
\setlength{\emergencystretch}{2em}
\multlinegap0pt
\usepackage{bm}
\usepackage{bbm}
\usepackage{dsfont}
\usepackage{xspace}
\usepackage{cancel}
\usepackage{mathtools}
\usepackage{amsthm}
\makeatletter
\@ifundefined{lemma}{\newtheorem{lemma}{Lemma}}{}
\makeatother
\newcommand{\hus}{\mathrel{{\leftarrow\mkern-11mu\ast}}}
\newcommand{\hut}{\mathrel{\leftarrow}}
\newcommand{\suh}{\mathrel{{\ast\mkern-11mu\to}}}
\newcommand{\sus}{\mathrel{{\ast\mkern-11mu\relbar\mkern-9mu\relbar\mkern-11mu\ast}}}
\newcommand{\tuh}{\mathrel{\rightarrow}}
\title[$\sigma$-Maximal Ancestral Graphs]{$\sigma$-Maximal Ancestral Graphs}
\author{Binghua Yao, Joris M. Mooij}
\date{Source: arXiv:2507.00093v1 (CC BY 4.0)}
\begin{document}
\maketitle

\noindent\emph{Source and licence.} $\sigma$-Maximal Ancestral Graphs. Version arXiv:2507.00093v1 (CC BY 4.0), distributed under the Creative Commons Attribution 4.0 International licence, as recorded in the arXiv licence metadata for this exact version and shown on the abstract page https://arxiv.org/abs/2507.00093v1. Full rights notice: licenses/arxiv-2507.00093-CC-BY-4.0.txt.\par\smallskip

\noindent\textit{Editorial scope.} Counted unit: the Lemma whose source label is \texttt{lemma3.3.10} in the article (source0.tex lines 1154--1161, source label \texttt{lemma3.3.10}), delivered together with the article's own complete original proof of it (source0.tex lines 1162--1182, one \texttt{proof} environment of 21 lines). No mathematics is written, repaired or re-derived: statement and proof are the article's own text line for line. Required context retained uncounted: none. Every definition, display and label of those lines is retained unchanged. Omitted from this package and not redistributed: 1--1153, 1183--2095. Nothing inside a retained excerpt is omitted or altered: every retained body is delivered exactly as the article's lines are written, so no omission receipt of any kind accompanies this package. Citations: the retained excerpts carry no citation command at all, so no bibliography entry of the article is transcribed and no reference is redistributed. Reference adaptation: none was needed and none was made; every reference inside the delivered unit resolves inside it, so no unresolved reference or citation is produced. Printed slips kept literal: no printed word, symbol, hypothesis or number of the counted statement or of its proof was altered. Layout and class: the article's own class is \texttt{uai2025} with packages including babel, natbib, mathtools, booktabs, tikz, amsthm, amssymb, algorithm, algorithmic, thm-restate, cleveref, subfig, thmtools; the wrapper is the lane's pinned \texttt{amsart} editorial wrapper at 10pt on A4, which restates the theorem-like environments the retained text needs and adds the article's own macro definitions that the retained text needs (\texttt{\textbackslash hus}, \texttt{\textbackslash hut}, \texttt{\textbackslash suh}, \texttt{\textbackslash sus}, \texttt{\textbackslash tuh}), with extra wrapper packages (bm, bbm, dsfont, xspace, cancel, mathtools). The delivered numbering is the unit's own. Assets: no retained span contains a figure environment, a graphics-inclusion command, a PDF-page inclusion, a table or a TikZ picture; no image file is copied into this package, the package declares no asset dependency and no image file is redistributed with it. Licence: version arXiv:2507.00093v1 of this article is distributed under the Creative Commons Attribution 4.0 International licence, as recorded in the arXiv licence metadata for this exact version and shown on the abstract page https://arxiv.org/abs/2507.00093v1. A full rights notice is delivered at licenses/arxiv-2507.00093-CC-BY-4.0.txt. Characters: every retained excerpt is pure ASCII, so the delivered engine, XeLaTeX with its default Latin Modern text font, reports no missing character.
\begin{lemma} \label{lemma3.3.10}
Let $\mathcal{G}=(\mathcal{V},\mathcal{E},\mathcal{L})$ be a DMG, $Z\subseteq \mathcal{V}$, and $\pi$ is a $Z$-$\sigma$-open walk between $v_0$ and $v_n$ in $\mathcal{G}$. Suppose $v_i\in\text{Sc}_\mathcal{G}(v_j)$ for some $i,j\in\{0,\ldots,n\}$ with $i<j$. If we then replace the subwalk $v_i\sus\cdots\sus v_j$ of $\pi$ by
\begin{enumerate}
    \item a shortest directed path $v_i\tuh\cdots\tuh v_j$ in $\mathcal{G}$ if $j=n$ or if $v_j\tuh v_{j+1}$ on $\pi$, or
    \item a shortest directed path $v_i\hut\cdots\hut v_j$ in $\mathcal{G}$ otherwise,
\end{enumerate}
then this new subwalk is entirely within $\text{Sc}_\mathcal{G}(v_j)$ and the modified walk $\pi'$ is still $Z$-$\sigma$-open.
\end{lemma}

\begin{proof}
$\pi^{\prime}$ cannot become $Z$-$\sigma$-blocked at one of the initial nodes $v_{0},\ldots,v_{i-1}$ or at one of the final nodes $v_{j+1},\ldots,v_{n}$ on $\pi^{\prime}$, since these nodes occur in the same local configuration on $\pi$ and are not $Z$-$\sigma$-blocked on $\pi$ by assumption. Furthermore, $\pi^{\prime}$ cannot become $Z$-$\sigma$-blocked at one of the nodes strictly between $v_{i}$ and $v_{j}$ on $\pi^{\prime}$ (if there are any), since these nodes are all non-endnode non-colliders that only point to nodes in the same strongly connected component $\mathrm{Sc}_\mathcal{G}(v_{j})$. It is also worth noting that $\pi^{\prime}$ cannot become $Z$-$\sigma$-blocked at any of its endnodes, which could be $v_{i}$ or $v_{j}$ or both, because those are the same in $\pi$. So in the following we can w.l.o.g. assume that both $v_{i}$ and $v_{j}$ are non-endnodes of $\pi$ and thus $\pi^{\prime}$.

Case 1: By assumption $v_{j}$ is in the subwalk $v_{j-1}\sus v_j\tuh v_{j+1}$ (or the right endnode) on $\pi$ that is $Z$-$\sigma$-open. Since the same blocking criteria apply to $v_{j}$ on $\pi^{\prime}$ it remains $Z$-$\sigma$-open on $\pi^{\prime}$. If $v_{i}=v_{j}$ then also $v_{i}$ is $Z$-$\sigma$-open on $\pi^{\prime}$ (if $v_{i}$ is the left endnode or not). If $v_{i}\neq v_{j}$, then the new directed path $v_{i}\tuh\cdots\tuh v_{j}$ in $\pi^{\prime}$ is $Z$-$\sigma$-open at $v_{i}$ because all nodes in between lie in the same strongly connected component $\mathrm{Sc}_\mathcal{G}(v_{i})$ (or $v_{i}$ is the left endnode anyways).

Case 2: Since case 1 is solved we can assume that we have $j<n$ with $v_{j}\hus v_{j+1}$ on $\pi$. If $v_{i-1}\hus v_{i}$ on $\pi^{\prime}$ (or $v_{i}$ the left endnode) then this case is analogous to case 1. So we can also assume that we have $i>0$ and $v_{i-1}\suh v_{i}$ on $\pi$. So $\pi$ looks as follows:
\[
\pi:\qquad\cdots v_{i-1}\suh v_{i}\sus\cdots \sus v_{j}\hus v_{j+1}\cdots.
\]
So there must be a smallest number $k\in\{i,\ldots,j\}$ such that a collider appears at $v_{k}$ on $\pi$:
\[
\pi:\qquad\cdots v_{i-1}\suh v_{i}\tuh\cdots \tuh v_{k}\hus\cdots\sus v_{j}\hus v_{j+1}\cdots.
\]
Since $\pi$ is $Z$-$\sigma$-open we have $v_{k}\in\mathrm{Anc}_\mathcal{G}(Z)$. Since $v_{i}\in\mathrm{Anc}_\mathcal{G}(v_{k})$ (otherwise $v_{k}$ would not be the first collider appearing after $v_{i}$) we thus have that also $v_{i}\in\mathrm{Anc}_\mathcal{G}(Z)$. So if we replace the subwalk $v_{i}\sus\cdots\sus v_{j}$ of $\pi$ by the shortest directed path $v_{i}\hut\cdots\hut v_{j}$ in $\mathcal{G}$ we then get for $\pi^{\prime}$ the following situation:
\[
\pi^{\prime}:\qquad\cdots v_{i-1}\suh v_{i}\hut\cdots \hut v_{j}\hus v_{j+1}\cdots,
\]
which is then $Z$-$\sigma$-open at $v_{i}$ as $v_{i}\in\mathrm{Anc}_\mathcal{G}(Z)$. Note that this holds also when $v_{i}=v_{j}$. If $v_{i}\neq v_{j}$ then $v_{j}$ is also $Z$-$\sigma$-open on $\pi^{\prime}$ as $v_{j}$ points left to a node in the same strongly connected component as $v_{j}$.

So in all cases $\pi^{\prime}$ stays $Z$-$\sigma$-open.
\end{proof}

\end{document}
